\documentclass[12pt,a4paper]{article}

% --- Metadata (per course) ---
\def\courseshortheader{THỐNG KÊ TOÁN}
\def\coverbookline{XÁC SUẤT VÀ THỐNG KÊ}
\def\covermaintitle{Thống kê toán}
\def\coversubtitle{Giáo án lý thuyết - Mẫu, ước lượng, khoảng tin cậy, kiểm định}

% Shared preamble for nhom_5 (requires \courseshortheader before \input)
\usepackage[utf8]{inputenc}
\usepackage[vietnamese]{babel}
\usepackage{amsmath,amssymb,amsthm}
\usepackage{geometry}
\usepackage{enumitem}
\usepackage[most]{tcolorbox}
\usepackage{hyperref}
\usepackage{fancyhdr}
\usepackage{graphicx}
\usepackage{booktabs}
\usepackage{tikz}
\usetikzlibrary{calc,arrows.meta,decorations.pathreplacing,positioning}
\usepackage{eso-pic}
\usepackage{xcolor}

\geometry{a4paper, margin=2cm, bottom=2.5cm}
\hypersetup{colorlinks=true, linkcolor=blue!70!black, urlcolor=blue}
\setlist[itemize]{topsep=4pt, itemsep=2pt}
\setlist[enumerate]{topsep=4pt, itemsep=2pt}

\AddToShipoutPictureFG{%
  \AtPageCenter{%
    \begin{tikzpicture}[remember picture, overlay]
      \node[rotate=45, scale=1, opacity=0.06]{%
        \fontsize{60}{72}\selectfont\bfseries\sffamily Đặng Tấn Quí%
      };
    \end{tikzpicture}%
  }%
}

\pagestyle{fancy}
\fancyhf{}
\fancyhead[L]{\small\textit{\courseshortheader}}
\fancyhead[R]{\small\textit{Đặng Tấn Quí}}
\fancyfoot[C]{\thepage}
\fancyfoot[L]{\footnotesize\textcopyright\ Đặng Tấn Quí}
\fancyfoot[R]{\footnotesize SĐT: 0377\,122\,210}
\renewcommand{\headrulewidth}{0.4pt}
\renewcommand{\footrulewidth}{0.4pt}

\fancypagestyle{plain}{%
  \fancyhf{}
  \fancyfoot[C]{\thepage}
  \fancyfoot[L]{\footnotesize\textcopyright\ Đặng Tấn Quí}
  \fancyfoot[R]{\footnotesize SĐT: 0377\,122\,210}
  \renewcommand{\headrulewidth}{0pt}
  \renewcommand{\footrulewidth}{0.4pt}
}

\newtcolorbox{dongco}[1][]{%
  enhanced, breakable,
  colback=teal!4!white, colframe=teal!60!black,
  fonttitle=\bfseries, title={Động cơ học -- Tại sao cần khái niệm này?}, #1%
}
\newtcolorbox{ynghia}[1][]{%
  enhanced, breakable,
  colback=purple!3!white, colframe=purple!50!black,
  fonttitle=\bfseries, title={Ý nghĩa \& Trực giác}, #1%
}
\newtcolorbox{ungdung}[1][]{%
  enhanced, breakable,
  colback=olive!5!white, colframe=olive!60!black,
  fonttitle=\bfseries, title={Ứng dụng thực tế}, #1%
}
\newtcolorbox{dinhnghia}[1][]{%
  enhanced, breakable,
  colback=blue!3!white, colframe=blue!60!black,
  fonttitle=\bfseries, title={Định nghĩa}, #1%
}
\newtcolorbox{dinhly}[1][]{%
  enhanced, breakable,
  colback=green!3!white, colframe=green!50!black,
  fonttitle=\bfseries, title={Định lý}, #1%
}
\newtcolorbox{tinhchat}[1][]{%
  enhanced, breakable,
  colback=yellow!5!white, colframe=orange!50!black,
  fonttitle=\bfseries, title={Tính chất}, #1%
}
\newtcolorbox{phuongphap}[1][]{%
  enhanced, breakable,
  colback=gray!5!white, colframe=gray!50!black,
  fonttitle=\bfseries, title={Phương pháp}, #1%
}
\newtcolorbox{luuy}[1][]{%
  enhanced, breakable,
  colback=red!3!white, colframe=red!50!black,
  fonttitle=\bfseries, title={Lưu ý}, #1%
}
\newtcolorbox{congthuc}[1][]{%
  enhanced, breakable,
  colback=cyan!3!white, colframe=cyan!50!black,
  fonttitle=\bfseries, title={Công thức}, #1%
}
\newtcolorbox{vidu}[1][]{%
  enhanced, breakable,
  colback=violet!3!white, colframe=violet!50!black,
  fonttitle=\bfseries, title={Ví dụ}, #1%
}
\newtcolorbox{phanvidu}[1][]{%
  enhanced, breakable,
  colback=red!2!white, colframe=red!40!black,
  fonttitle=\bfseries, title={Phản ví dụ}, #1%
}
\newtcolorbox{tongket}[1][]{%
  enhanced, breakable,
  colback=blue!4!white, colframe=blue!70!black,
  fonttitle=\bfseries, title={Tổng kết chương}, #1%
}


\begin{document}

\thispagestyle{empty}
\begin{tikzpicture}[remember picture, overlay]
  \fill[blue!80!black] (current page.south west) rectangle (current page.north east);
  \fill[blue!60!cyan, opacity=0.8]
    (current page.south west) -- (current page.north west) --
    ([xshift=-3cm]current page.north east) -- (current page.south east) -- cycle;
  \foreach \r/\o in {6/0.05, 4.5/0.08, 3/0.12} {
    \fill[white, opacity=\o] ([xshift=2cm, yshift=-2cm]current page.north east) circle (\r cm);
  }
  \draw[white, line width=2pt] ([yshift=-5cm]current page.north west) ++(2cm,0) -- ++(17cm,0);
  \draw[white, line width=2pt] ([yshift=-16cm]current page.north west) ++(2cm,0) -- ++(17cm,0);
  \node[anchor=west, white, font=\fontsize{28}{34}\bfseries\selectfont]
    at ([xshift=3cm, yshift=-7.5cm]current page.north west) {\coverbookline};
  \node[anchor=west, white, font=\fontsize{20}{26}\selectfont]
    at ([xshift=3cm, yshift=-9cm]current page.north west) {\covermaintitle};
  \node[anchor=west, white, font=\fontsize{16}{22}\selectfont]
    at ([xshift=3cm, yshift=-10.2cm]current page.north west) {\coversubtitle};
  \node[anchor=west, white, font=\Large\bfseries]
    at ([xshift=3cm, yshift=-13cm]current page.north west) {Biên soạn:};
  \node[anchor=west, yellow!80!white, font=\fontsize{24}{30}\bfseries\selectfont]
    at ([xshift=3cm, yshift=-14.5cm]current page.north west) {ĐẶNG TẤN QUÍ};
  \node[anchor=west, white!90, font=\large]
    at ([xshift=3cm, yshift=-18cm]current page.north west) {SĐT: 0377\,122\,210};
  \node[anchor=south, white!80, font=\small]
    at ([yshift=1.5cm]current page.south) {Tài liệu có bản quyền --- Cấm sao chép khi chưa được sự đồng ý của tác giả};
  \node[anchor=south east, white, font=\Large\bfseries]
    at ([xshift=-2cm, yshift=3cm]current page.south east) {\the\year};
\end{tikzpicture}

\newpage
\tableofcontents
\newpage

\section*{Lời nói đầu}
\addcontentsline{toc}{section}{Lời nói đầu}

Tài liệu thuộc \textbf{Nhóm 5 --- Xác suất và thống kê}, biên soạn theo cùng triết lý với giáo án Đại số tuyến tính: học để \emph{hiểu} và \emph{dùng}, không chỉ để ghi nhớ công thức.

\begin{enumerate}
  \item \textbf{Động cơ trước định nghĩa.} Mỗi phần lớn mở bằng ô \textsf{\color{teal!60!black} Động cơ học} --- câu hỏi thực tế hoặc bài toán mẫu cần khái niệm đó.
  \item \textbf{Trực giác trước công thức.} Ô \textsf{\color{purple!50!black} Ý nghĩa \& Trực giác} giải thích ``công thức đang nói gì'' (xác suất = tần suất dài hạn, kỳ vọng = trung bình có trọng số, v.v.).
  \item \textbf{Ví dụ có kiểm tra.} Ô \textsf{\color{violet!50!black} Ví dụ} nên tự làm trước khi đọc lời giải; phần thống kê ứng dụng cần gắn dữ liệu số cụ thể.
  \item \textbf{Tổng kết cuối mỗi chương.} Ô \textsf{\color{blue!70!black} Tổng kết chương} liệt kê kết quả phải nhớ và liên kết sang phần sau.
  \item \textbf{Phân tầng theo môn.} X1--X3 là nền; X4--X5 nâng cao lý thuyết; X6--X8 chuyên sâu (đa biến, Bayes, quá trình ngẫu nhiên).
\end{enumerate}

\noindent\textbf{Cách đọc:} Động cơ $\to$ Định nghĩa / Định lý $\to$ Ví dụ $\to$ Ý nghĩa $\to$ Tổng kết. Với môn thống kê, luôn hỏi: \emph{mẫu là gì, tham số nào chưa biết, giả thuyết $H_0$ là gì}.

\newpage


\section*{Tại sao học thống kê toán?}

\begin{dongco}[title={Động cơ học - Ba tình huống quen thuộc}]
\textbf{Câu chuyện 1 - Tỷ lệ phế phẩm nhà máy.}
Một nhà máy sản xuất 5\,000\,000 sản phẩm; không thể kiểm tra hết. Chọn 500 sản phẩm, phát hiện 20 lỗi $\Rightarrow$ tần suất mẫu $\dfrac{20}{500} = 0{,}04 = 4\%$. Con số này \emph{gợi ý} tỷ lệ phế phẩm tổng thể, nhưng cần ngôn ngữ xác suất để đo ``độ không chắc chắn''.

\medskip
\textbf{Câu chuyện 2 - Trọng lượng trung bình.}
Cân 36 sản phẩm, trung bình mẫu $\bar{x} = 166{,}22$ gam. Điểm ước lượng chưa đủ: ta cần khoảng tin cậy $[164{,}26;\;168{,}18]$ với độ tin cậy 95\%.

\medskip
\textbf{Câu chuyện 3 - Kiểm định máy sản xuất.}
Máy quy định trung bình 100 gam; cân 100 sản phẩm được $\bar{x} = 100{,}4$. Liệu đây chỉ là dao động ngẫu nhiên hay máy đã lệch? Kiểm định giả thuyết trả lời câu hỏi đó với mức ý nghĩa $\alpha = 5\%$.

\medskip
Ba câu chuyện trên cùng một khung: \textbf{mẫu ngẫu nhiên} $\to$ \textbf{thống kê mẫu} $\to$ \textbf{phân phối mẫu} $\to$ \textbf{ước lượng} $\to$ \textbf{khoảng tin cậy} $\to$ \textbf{kiểm định}.
\end{dongco}

\newpage

\section{Thống kê và ước lượng tham số}
\subsection{Mẫu ngẫu nhiên và phân phối mẫu}

\subsubsection{Tổng thể và mẫu}

\begin{dongco}[title={Động cơ - Tại sao lấy mẫu?}]
Khảo sát dấu hiệu $X$ trên tập phần tử (tổng thể) thường không khả thi vì: quy mô quá lớn, chi phí/thời gian, hoặc phép đo phá hủy đối tượng. Thay vì khảo sát toàn bộ, ta \textbf{lấy mẫu} - tập con có cỡ $n$ - rồi suy luận về tổng thể.
\end{dongco}

\begin{dinhnghia}[title={Tổng thể và mẫu}]
  \textbf{Tổng thể}: tập tất cả phần tử cần nghiên cứu; ký hiệu $N$ là số phần tử (hữu hạn hoặc vô hạn).
  
  \textbf{Dấu hiệu} $X$ là đại lượng định lượng hoặc định tính trên mỗi phần tử.

  \textbf{Phép lấy mẫu}: chọn tập con từ tổng thể. \textbf{Mẫu} là tập con đó; \textbf{cỡ mẫu} (kích thước mẫu) ký hiệu $n$.
\end{dinhnghia}

\begin{vidu}[title={Tổng thể - thu nhập và phế phẩm}]
\begin{itemize}
  \item Điều tra thu nhập hộ gia đình Hà Nội: tổng thể = mọi hộ; dấu hiệu = thu nhập (định lượng).
  \item Nhà máy 5\,000\,000 sản phẩm, đánh giá tỷ lệ phế phẩm: tổng thể = 5\,000\,000 sản phẩm; dấu hiệu = ``có lỗi hay không'' (định tính $\to$ Bernoulli).
\end{itemize}
\end{vidu}

\begin{vidu}[title={Mẫu - tỷ lệ phế phẩm nhà máy}]
Không kiểm tra hết 5\,000\,000 sản phẩm; chọn $n = 500$, phát hiện 20 lỗi.
\[
  \hat{p} = \dfrac{20}{500} = 0{,}04 = 4\%.
\]
\end{vidu}

\begin{luuy}[title={Mẫu phải đại diện}]
Nếu chỉ khảo sát sinh viên ngành CNTT để ước lượng lương trung bình SV mới ra trường, kết quả \textbf{chệch mẫu}. Cần lấy mẫu ngẫu nhiên, phân tầng, hoặc theo khối đúng quy trình.
\end{luuy}

\subsubsection{Biểu diễn dữ liệu}

\begin{dinhnghia}[title={Các dạng bảng dữ liệu}]
  Từ mẫu cỡ $n$, ghi $(x_1, x_2, \ldots, x_n)$.
  \begin{itemize}
    \item \textbf{Dạng liệt kê:} liệt kê từng giá trị.
    \item \textbf{Dạng rút gọn:} bảng tần số $n_i$ hoặc tần suất $f_i = \dfrac{n_i}{n}$ tại các giá trị $x_i$.
    \item \textbf{Dạng khoảng:} chia $(a,b)$ thành các lớp $(a_{i-1}, a_i]$; dùng điểm giữa $x_i = \tfrac{1}{2}(a_{i-1}+a_i)$ khi tính toán.
  \end{itemize}
\end{dinhnghia}

\begin{vidu}[title={Biểu đồ - xếp hạng sản phẩm}]
Khảo sát 400 khách hàng: A (35), B (260), C (93), D (12). Tần suất: 8{,}75\%; 65\%; 23{,}25\%; 3\%. Biểu đồ cột/quạt giúp so sánh trực quan nhóm định tính.
\end{vidu}

\begin{congthuc}[title={Biểu đồ tần suất}]
  \textbf{Tần suất tích lũy:} $F_n(x) = \dfrac{1}{n}\displaystyle\sum_{i=1}^n \mathbf{1}(x_i \le x)$.
  
  \textbf{Biểu đồ tần suất:} diện tích mỗi cột $\propto$ tần suất lớp; chiều cao $= \dfrac{\text{tần suất}}{\text{độ rộng lớp}}$.
\end{congthuc}

\subsubsection{Các đặc trưng mô tả}

\begin{congthuc}[title={Vị trí trung tâm - mẫu quan sát}]
  \textbf{Trung bình mẫu:}
  \[
    \bar{x} = \dfrac{1}{n}\displaystyle\sum_{i=1}^n x_i
    \quad\text{hoặc}\quad
    \bar{x} = \dfrac{1}{n}\displaystyle\sum_{i=1}^k n_i x_i \;\;(\text{dạng rút gọn}).
  \]
  \textbf{Trung vị mẫu} $\tilde{x}$: sắp $x_{(1)} \le \cdots \le x_{(n)}$;
  \[
    \tilde{x} =
    \begin{cases}
      x_{(\frac{n+1}{2})}, & n \text{ lẻ},\\[4pt]
      \tfrac{1}{2}(x_{(\frac{n}{2})} + x_{(\frac{n}{2}+1)}), & n \text{ chẵn}.
    \end{cases}
  \]
  \textbf{Mốt:} giá trị có tần số lớn nhất (có thể không duy nhất).
\end{congthuc}

\begin{congthuc}[title={Phân tán - phương sai và độ lệch chuẩn mẫu}]
  \textbf{Phương sai mẫu} (chưa hiệu chỉnh):
  \[
    s_e^2 = \dfrac{1}{n}\displaystyle\sum_{i=1}^n (x_i - \bar{x})^2
    \quad\text{hoặc}\quad
    s_e^2 = \dfrac{1}{n}\displaystyle\sum_{i=1}^n n_i(x_i - \bar{x})^2 \;\;(\text{dạng rút gọn}).
  \]
  \textbf{Phương sai mẫu hiệu chỉnh} (dùng trong suy luận):
  \[
    s^2 = \dfrac{1}{n-1}\displaystyle\sum_{i=1}^n (x_i - \bar{x})^2, \qquad s = \sqrt{s^2}.
  \]
  Dạng tính nhanh: 
  \[
    s_e^2 = \dfrac{1}{n}\displaystyle\sum x_i^2 - \left(\dfrac{1}{n}\displaystyle\sum x_i\right)^2
    \quad\text{hoặc}\quad
    s_e^2 = \dfrac{1}{n}\displaystyle\sum n_i x_i^2 - \left(\dfrac{1}{n}\displaystyle\sum n_i x_i\right)^2 \;\;(\text{dạng rút gọn}).
  \]
\end{congthuc}

\begin{congthuc}[title={Độ lệch, độ nhọn, tứ phân vị}]
\begin{itemize}
  \item \textbf{Độ lệch mẫu:} $g_1 = \dfrac{m_3}{s^3}$, $m_3 = \dfrac{1}{n}\displaystyle\sum(x_i-\bar{x})^3$.
  \item \textbf{Độ nhọn:} $g_2 = \dfrac{m_4}{s^4} - 3$.
  \item \textbf{Tứ phân vị:} $Q_1 = x_{(0{,}25n)}$, $Q_3 = x_{(0{,}75n)}$.
  \item \textbf{Khoảng tứ phân vị:} $IQR = Q_3 - Q_1$.
\end{itemize}
\end{congthuc}

\begin{vidu}[title={Số lỗi trên sản phẩm đúc}]
150 sản phẩm, bảng tần số theo số lỗi 0-6. Kết quả:
\[
  \bar{x} = \dfrac{499}{150} \approx 3{,}3267, \qquad
  s_e^2 \approx 2{,}7533, \qquad
  s_e \approx 1{,}6593.
\]
\end{vidu}

\begin{vidu}[title={Phương sai mẫu hiệu chỉnh - giá tăng}]
Bốn siêu thị, mức tăng giá (nghìn đồng/kg): 12, 15, 17, 20.
\[
  \bar{x} = 16, \qquad s^2 = \dfrac{(12-16)^2 + (15-16)^2 + (17-16)^2 + (20-16)^2}{3} = \dfrac{34}{3}.
\]
\end{vidu}

\subsubsection{Mẫu ngẫu nhiên}

\begin{dinhnghia}[title={Mẫu ngẫu nhiên}]
  Cho biến ngẫu nhiên gốc $X$ có phân phối $F_X(x)$. Bộ $(X_1, X_2, \ldots, X_n)$ là \textbf{mẫu ngẫu nhiên} cỡ $n$, ký hiệu $W_X = (X_1, \ldots, X_n)$, nếu $X_1, \ldots, X_n$ \textbf{độc lập} và \textbf{cùng phân phối} $F_X(x)$.

  Hàm phân phối đồng thời:
  \[
    F_{X_1,\ldots,X_n}(x_1,\ldots,x_n) = \prod_{i=1}^n F_X(x_i).
  \]
  \textbf{Mẫu cụ thể} (quan sát): $W_x = (x_1, \ldots, x_n)$.
\end{dinhnghia}

\begin{vidu}[title={Pin loại A}]
Tuổi thọ pin $X \sim \mathcal{N}(\mu, \sigma^2)$. Chọn $n=10$ pin, $X_i$ = tuổi thọ pin thứ $i$. Một quan sát: $W_x = (1000; 1001; \ldots; 1009)$.
\end{vidu}

\begin{dinhnghia}[title={Thống kê}]
  \textbf{Thống kê} là hàm $T = g(X_1, \ldots, X_n)$ \emph{không} chứa tham số tổng thể chưa biết. Thống kê là biến ngẫu nhiên; giá trị tại mẫu cụ thể gọi là \textbf{giá trị quan sát}.
\end{dinhnghia}

\begin{congthuc}[title={Thống kê mẫu thông dụng}]
  \begin{itemize}
    \item \textbf{Trung bình mẫu ngẫu nhiên:} $\bar{X} = \dfrac{1}{n}\displaystyle\sum_{i=1}^n X_i$.
    \item \textbf{Phương sai mẫu hiệu chỉnh:} $S^2 = \dfrac{1}{n-1}\displaystyle\sum_{i=1}^n (X_i - \bar{X})^2$.
    \item \textbf{Tần suất mẫu} (Bernoulli): $\hat{P} = \dfrac{X}{n}$ với $X = \displaystyle\sum X_i$, $X_i \in \{0,1\}$.
  \end{itemize}
\end{congthuc}

\begin{tinhchat}[title={Kỳ vọng và phương sai của $\bar{X}$, $\hat{P}$}]
  Nếu $E(X) = \mu$, $V(X) = \sigma^2$:
  \[
    E(\bar{X}) = \mu, \qquad V(\bar{X}) = \dfrac{\sigma^2}{n}.
  \]
  Nếu $X_i \sim \text{Bernoulli}(p)$:
  \[
    E(\hat{P}) = p, \qquad V(\hat{P}) = \dfrac{p(1-p)}{n}.
  \]
  $E(S^2) = \sigma^2$ (hiệu chỉnh $n-1$); còn $E(S_e^2) = \dfrac{n-1}{n}\sigma^2$ (có chệch).
\end{tinhchat}

\begin{ynghia}[title={Trực giác $V(\bar{X}) = \dfrac{\sigma^2}{n}$}]
Trung bình mẫu ``ổn định'' hơn một quan sát đơn: phương sai giảm $n$ lần. Cỡ mẫu lớn $\Rightarrow$ ước lượng $\bar{X}$ tập trung quanh $\mu$ hơn - nền tảng của CLT và khoảng tin cậy hẹp hơn.
\end{ynghia}

\subsubsection{Phân phối mẫu}

\begin{dinhnghia}[title={Phân phối mẫu}]
  Vì thống kê $T = g(X_1,\ldots,X_n)$ là biến ngẫu nhiên, \textbf{phân phối xác suất của $T$} gọi là \textbf{phân phối mẫu} của thống kê đó.
\end{dinhnghia}

\begin{dinhly}[title={Phân phối mẫu của $\bar{X}$ - tổng thể chuẩn}]
  Nếu $X \sim \mathcal{N}(\mu, \sigma^2)$ thì
  \[
    \bar{X} \sim \mathcal{N}\!\left(\mu,\; \dfrac{\sigma^2}{n}\right).
  \]
\end{dinhly}

\begin{vidu}[title={Điện trở - xác suất trung bình mẫu $< 95\,\Omega$}]
Điện trở $X \sim \mathcal{N}(100, 8^2)$. Mẫu $n=16$:
\[
  \bar{X} \sim \mathcal{N}\!\left(100,\; \dfrac{64}{16}\right) = \mathcal{N}(100, 2^2).
\]
\[
  P(\bar{X} < 95) = \Phi\!\left(\dfrac{95-100}{2}\right) = \Phi(-2{,}5) \approx 0{,}00621.
\]
\end{vidu}

\begin{dinhly}[title={Định lý giới hạn trung tâm (CLT)}]
  Mẫu gồm các $X_i$ \textbf{độc lập, cùng phân phối} (lấy từ $X$) với $E(X)=\mu$, $V(X)=\sigma^2 < \infty$. Khi $n \to \infty$,
  \[
    Z = \dfrac{\bar{X} - \mu}{\sigma/\sqrt{n}} \xrightarrow{d} \mathcal{N}(0,1).
  \]
  Thực hành: $n \ge 30$ thì $Z$ \textbf{xấp xỉ} $\mathcal{N}(0,1)$ dù $X$ không chuẩn.
\end{dinhly}

\begin{vidu}[title={CLT - phân phối đều}]
$X \sim \mathcal{U}[2,4]$: $\mu = 3$, $\sigma^2 = 1/3$. Với $n=40$,
\[
  \bar{X} \approx \mathcal{N}\!\left(3,\; \dfrac{1/3}{40}\right) = \mathcal{N}\!\left(3,\; \dfrac{1}{120}\right).
\]
\end{vidu}

\begin{dinhly}[title={Phân phối $\chi^2$, $t$ (Student), $F$ (Fisher)}]
  Với $X \sim \mathcal{N}(\mu, \sigma^2)$, mẫu cỡ $n$:
  \begin{itemize}
    \item $\dfrac{(n-1)S^2}{\sigma^2} \sim \chi^2(n-1)$.
    \item $\dfrac{\bar{X}-\mu}{S/\sqrt{n}} \sim t(n-1)$.
    \item $\bar{X}$ và $S^2$ \textbf{độc lập}.
    \item Hai mẫu chuẩn độc lập: $F = \dfrac{S_1^2/\sigma_1^2}{S_2^2/\sigma_2^2} \sim F(n_1-1, n_2-1)$.
  \end{itemize}
\end{dinhly}

\begin{congthuc}[title={Phân phối mẫu $\bar{X}_1 - \bar{X}_2$}]
  Hai tổng thể độc lập, $\bar{X}_1$, $\bar{X}_2$ từ mẫu cỡ $n_1$, $n_2$:
  \[
    E(\bar{X}_1 - \bar{X}_2) = \mu_1 - \mu_2, \qquad
    V(\bar{X}_1 - \bar{X}_2) = \dfrac{\sigma_1^2}{n_1} + \dfrac{\sigma_2^2}{n_2}.
  \]
  Nếu $n_1, n_2 \ge 30$ (hoặc chuẩn): chuẩn hóa $\approx \mathcal{N}(0,1)$.
\end{congthuc}

\begin{luuy}[title={Khi nào dùng phân phối nào?}]
\begin{itemize}
  \item $\sigma$ biết, $X$ chuẩn: dùng $Z = \dfrac{\bar{X}-\mu}{\sigma/\sqrt{n}}$.
  \item $\sigma$ chưa biết, $X$ chuẩn, $n$ nhỏ: dùng $t(n-1)$.
  \item $n$ lớn, không chắc chuẩn: CLT $\Rightarrow$ dùng $Z$ với $S$ thay $\sigma$.
\end{itemize}
\end{luuy}

\begin{tongket}[title={Tổng kết mục Mẫu ngẫu nhiên và phân phối mẫu}]
\begin{enumerate}
  \item Phân biệt \textbf{tổng thể} ($N$) và \textbf{mẫu} ($n$); mẫu phải đại diện.
  \item Thuộc công thức $\bar{x}$, $s^2$, $\tilde{x}$; biết $s_e^2$ vs $s^2$.
  \item Mẫu ngẫu nhiên = các quan sát độc lập, cùng phân phối; thống kê = hàm của mẫu.
  \item Phân phối mẫu: $\bar{X}$ chuẩn nếu $X$ chuẩn; CLT khi $n$ lớn; $\chi^2$, $t$, $F$ cho suy luận.
\end{enumerate}
\end{tongket}
\subsection{Ước lượng điểm}

\subsubsection{Khái niệm ước lượng điểm}

\begin{dongco}[title={Động cơ - Từ mẫu suy ra tham số}]
Tổng thể có tham số $\theta$ (ví dụ $\mu$, $\sigma^2$, $p$) chưa biết. Từ mẫu ngẫu nhiên $W_X = (X_1,\ldots,X_n)$ ta xây hàm $\hat{\Theta} = g(X_1,\ldots,X_n)$ - \textbf{ước lượng điểm} cho $\theta$. Giá trị quan sát $\hat{\theta} = g(x_1,\ldots,x_n)$ là con số cụ thể.
\end{dongco}

\begin{dinhnghia}[title={Ước lượng điểm}]
  Cho tham số $\theta$ của biến ngẫu nhiên $X$. \textbf{Ước lượng điểm} của $\theta$ là thống kê $\hat{\Theta} = g(X_1,\ldots,X_n)$.
  Giá trị $\hat{\theta} = g(x_1,\ldots,x_n)$ gọi là \textbf{giá trị ước lượng điểm}.
\end{dinhnghia}

\begin{congthuc}[title={Các ước lượng điểm thông dụng}]
  \begin{itemize}
    \item Kỳ vọng: $\hat{\mu} = \bar{X}$ (trung bình mẫu).
    \item Phương sai: $\hat{\sigma}^2 = S^2$ (phương sai mẫu hiệu chỉnh).
    \item Tỷ lệ: $\hat{p} = \dfrac{X}{n}$ (tần suất mẫu), $X$ = số phần tử có dấu hiệu A.
    \item Hiệu hai kỳ vọng: $\widehat{\mu_1-\mu_2} = \bar{X}_1 - \bar{X}_2$.
    \item Hiệu hai tỷ lệ: $\hat{p}_1 - \hat{p}_2$.
  \end{itemize}
\end{congthuc}

\begin{vidu}[title={Ước lượng $\mu$ và $\sigma^2$}]
$X \sim \mathcal{N}(\mu, \sigma^2)$, mẫu $n=10$: 13{,}8; 10{,}4; 9{,}7; 12{,}6; 14{,}1; 10{,}8; 15{,}1; 9{,}5; 13{,}1; 11{,}3.
\[
  \hat{\mu} = \bar{x} = 12{,}04, \qquad \hat{\sigma}^2 = s^2 \approx 4{,}8783.
\]
\end{vidu}

\begin{vidu}[title={Nhiều cách ước lượng $\mu$}]
Cùng dữ liệu trên, có thể chọn:
\begin{itemize}
  \item Trung bình mẫu: $\bar{x} = 12{,}04$.
  \item Trung vị: $\tilde{x} = \dfrac{11{,}3 + 12{,}6}{2} = 11{,}95$.
  \item Trung bình sau bỏ min/max: $\bar{x}_{-(\min+\max)} = 11{,}975$.
\end{itemize}
Cần tiêu chuẩn để chọn hàm ước lượng ``tốt''.
\end{vidu}

\subsubsection{Tiêu chuẩn lựa chọn ước lượng}

\begin{dinhnghia}[title={Ước lượng không chệch}]
  $\hat{\Theta}$ là \textbf{ước lượng không chệch} cho $\theta$ nếu $E(\hat{\Theta}) = \theta$ với mọi $\theta$.
\end{dinhnghia}

\begin{tinhchat}[title={Không chệch của $\bar{X}$, $S^2$, $\hat{P}$}]
  $E(\bar{X}) = \mu$; $E(S^2) = \sigma^2$; $E(\hat{P}) = p$.
  
  Ngược lại $E(S_e^2) = \dfrac{n-1}{n}\sigma^2$ nên $S_e^2$ \textbf{chệch}.
\end{tinhchat}

\begin{dinhnghia}[title={Ước lượng vững}]
  $\hat{\Theta}_n$ \textbf{vững} cho $\theta$ nếu $\hat{\Theta}_n \xrightarrow{P} \theta$ khi $n \to \infty$.
  
  Trực giác: cỡ mẫu lớn thì ước lượng ``sát'' tham số thật.
\end{dinhnghia}

\begin{tinhchat}[title={Vững của trung bình mẫu (ĐL lớn số)}]
  Với $X_i$ độc lập, cùng phân phối, $E(X)=\mu$, $V(X)=\sigma^2<\infty$: $\bar{X} \xrightarrow{P} \mu$. Do đó $\bar{X}$ vững cho $\mu$.
\end{tinhchat}

\begin{congthuc}[title={Phương sai và sai số tiêu chuẩn của ước lượng}]
  \textbf{Phương sai ước lượng:} $V(\hat{\Theta})$.
  
  \textbf{Sai số tiêu chuẩn:} $\text{SE}(\hat{\Theta}) = \sqrt{V(\hat{\Theta})}$.
  
  Ví dụ: $V(\bar{X}) = \sigma^2/n$, $\text{SE}(\bar{X}) = \sigma/\sqrt{n}$.
\end{congthuc}

\begin{dinhnghia}[title={Sai số bình phương trung bình - MSE}]
  \[
    \text{MSE}(\hat{\Theta}) = E\bigl[(\hat{\Theta}-\theta)^2\bigr] = V(\hat{\Theta}) + \bigl(E(\hat{\Theta})-\theta\bigr)^2.
  \]
  MSE = phương sai + bình phương độ chệch. Ước lượng không chệch: MSE = $V(\hat{\Theta})$.
\end{dinhnghia}

\begin{ynghia}[title={Đánh đổi chệch-phương sai}]
Đôi khi ước lượng hơi chệch nhưng phương sai nhỏ hơn nhiều $\Rightarrow$ MSE nhỏ hơn ước lượng không chệch (ví dụ hồi quy ridge trong học máy). Ưu tiên không chệch + phương sai nhỏ.
\end{ynghia}

\begin{luuy}[title={So sánh hai ước lượng không chệch}]
Nếu $\hat{\Theta}_1$, $\hat{\Theta}_2$ đều không chệch, chọn cái có $V(\hat{\Theta})$ nhỏ hơn (hiệu quả hơn relative).
\end{luuy}

\subsubsection{Phương pháp mô men}

\begin{phuongphap}[title={Phương pháp mô men (MME)}]
\begin{itemize}
  \item \textbf{Bước 1:} Tính mô men lý thuyết $\mu_k(\theta) = E(X^k)$.
  \item \textbf{Bước 2:} Tính mô men mẫu $\hat{\mu}_k = \dfrac{1}{n}\displaystyle\sum X_i^k$.
  \item \textbf{Bước 3:} Giải hệ $\mu_k(\theta) = \hat{\mu}_k$ ($k=1,2,\ldots$) để có $\hat{\theta}$.
\end{itemize}
\end{phuongphap}

\begin{vidu}[title={MME cho phân phối Poisson}]
$X \sim \text{Poisson}(\lambda)$: $E(X) = \lambda$. Đặt $\bar{X} = \lambda$ $\Rightarrow$ $\hat{\lambda}_{\text{MME}} = \bar{X}$.
\end{vidu}

\begin{vidu}[title={MME cho phân phối Uniform$(0,\theta)$}]
$E(X) = \theta/2$. Giải $\bar{X} = \theta/2$ $\Rightarrow$ $\hat{\theta}_{\text{MME}} = 2\bar{X}$.
\end{vidu}

\begin{vidu}[title={MME cho phân phối chuẩn}]
$E(X) = \mu$, $V(X) = \sigma^2$. Giải:
\[
  \hat{\mu} = \bar{X}, \qquad \hat{\sigma}^2 = \dfrac{1}{n}\displaystyle\sum(X_i-\bar{X})^2 = S_e^2.
\]
Lưu ý: $\hat{\sigma}^2$ theo MME \textbf{chệch}; phiên bản hiệu chỉnh $S^2$ thường được dùng hơn.
\end{vidu}

\subsubsection{Phương pháp hợp lý cực đại (MLE)}

\begin{phuongphap}[title={MLE}]
\begin{itemize}
  \item \textbf{Hàm hợp lý:} $L(\theta) = \prod_{i=1}^n f(x_i;\theta)$.
  \item \textbf{Log-hợp lý} (nhật ký hàm hợp lý): $\ell(\theta) = \ln L(\theta) = \displaystyle\sum_{i=1}^n \ln f(x_i;\theta)$.
  \item \textbf{MLE:} $\hat{\theta}_{\text{MLE}} = \arg\max_\theta L(\theta)$, thường giải $\dfrac{\partial \ell}{\partial \theta} = 0$.
\end{itemize}
\end{phuongphap}

\begin{vidu}[title={MLE cho Bernoulli$(p)$}]
$L(p) = p^{\displaystyle\sum x_i}(1-p)^{n-\displaystyle\sum x_i}$. $\ell(p) = (\displaystyle\sum x_i)\ln p + (n-\displaystyle\sum x_i)\ln(1-p)$.
\[
  \dfrac{\partial \ell}{\partial p} = 0 \;\Rightarrow\; \hat{p}_{\text{MLE}} = \dfrac{\displaystyle\sum x_i}{n} = \hat{p}.
\]
\end{vidu}

\begin{vidu}[title={MLE cho $\mathcal{N}(\mu,\sigma^2)$}]
Giải $\dfrac{\partial\ell}{\partial\mu} = 0$, $\dfrac{\partial\ell}{\partial\sigma^2} = 0$:
\[
  \hat{\mu}_{\text{MLE}} = \bar{X}, \qquad \hat{\sigma}^2_{\text{MLE}} = \dfrac{1}{n}\displaystyle\sum(X_i-\bar{X})^2 = S_e^2.
\]
\end{vidu}

\begin{vidu}[title={MLE cho Poisson$(\lambda)$}]
$\hat{\lambda}_{\text{MLE}} = \bar{X}$ (trùng MME).
\end{vidu}

\begin{vidu}[title={MLE cho $\mathcal{U}(0,\theta)$}]
Hàm hợp lý $L(\theta) = \theta^{-n}$ nếu $\theta \ge \max x_i$, bằng 0 nếu không. $\Rightarrow$ $\hat{\theta}_{\text{MLE}} = \max(X_1,\ldots,X_n)$.
\end{vidu}

\begin{tinhchat}[title={Tính chất MLE (tóm tắt)}]
  \begin{itemize}
    \item \textbf{Bất biến:} nếu $\hat{\theta}$ là MLE thì $g(\hat{\theta})$ là MLE cho $g(\theta)$.
    \item \textbf{Tiệm cận không chệch} và \textbf{tiệm cận hiệu quả} (đạt cận Cramér--Rao).
    \item \textbf{Tiệm cận chuẩn:} $\sqrt{n}(\hat{\theta}-\theta) \xrightarrow{d} \mathcal{N}(0, I(\theta)^{-1})$.
  \end{itemize}
\end{tinhchat}

\begin{congthuc}[title={Thông tin Fisher}]
  \[
    I(\theta) = E\!\left[\left(\dfrac{\partial}{\partial\theta}\ln f(X;\theta)\right)^2\right] = -E\!\left[\dfrac{\partial^2}{\partial\theta^2}\ln f(X;\theta)\right].
  \]
  Với mẫu độc lập, cùng phân phối, cỡ $n$: $I_n(\theta) = n\,I(\theta)$.
\end{congthuc}

\begin{luuy}[title={MME và MLE}]
MME dễ tính, trực quan; MLE thường hiệu quả hơn (phương sai nhỏ hơn tiệm cận). Với nhiều mô hình chuẩn, hai phương pháp cho cùng kết quả (Bernoulli, Poisson, $\mu$ chuẩn).
\end{luuy}

\begin{tongket}[title={Tổng kết}]
\begin{enumerate}
  \item $\hat{\Theta} = g(X_1,\ldots,X_n)$; giá trị quan sát $\hat{\theta}$.
  \item Không chệch: $E(\hat{\Theta})=\theta$; vững: $\hat{\Theta}_n \xrightarrow{P}\theta$; MSE = $V$ + chệch$^2$.
  \item MME: đồng bộ mô men lý thuyết và mô men mẫu.
  \item MLE: cực đại $L(\theta)$; thường hiệu quả tiệm cận.
  \item Thuộc ước lượng thông dụng: $\bar{X}$, $S^2$, $\hat{p}$.
\end{enumerate}
\end{tongket}
\subsection{Ước lượng khoảng}

\subsubsection{Khái niệm khoảng tin cậy}

\begin{dongco}[title={Động cơ - Điểm ước lượng chưa đủ}]
$\bar{x} = 166{,}22$ gam cho biết ``điểm tốt nhất'' cho $\mu$, nhưng mẫu khác sẽ cho $\bar{x}$ khác. Ta cần \textbf{khoảng} bao quanh $\mu$ với xác suất đã cho - đó là khoảng tin cậy.
\end{dongco}

\begin{dinhnghia}[title={Khoảng tin cậy (KTC)}]
  Cho tham số $\theta$ và độ tin cậy $1-\alpha \in (0,1)$. Khoảng ngẫu nhiên $[\hat{\Theta}_L, \hat{\Theta}_U]$ là \textbf{khoảng tin cậy} mức $1-\alpha$ cho $\theta$ nếu
  \[
    P(\hat{\Theta}_L \le \theta \le \hat{\Theta}_U) \ge 1 - \alpha.
  \]
  \textbf{Lưu ý:} $\theta$ là hằng số; ngẫu nhiên là \emph{cận} $[\hat{\Theta}_L, \hat{\Theta}_U]$, không phải $\theta$.
\end{dinhnghia}

\begin{luuy}[title={Diễn giải đúng độ tin cậy 95\%}]
Nếu lặp lại lấy mẫu rất nhiều lần và mỗi lần tính KTC 95\%, thì khoảng \textbf{95\% số lần} sẽ chứa $\theta$ thật. \textbf{Không} viết $P(164{,}26 \le \mu \le 168{,}18) = 0{,}95$ cho một mẫu cụ thể.
\end{luuy}

\begin{ynghia}[title={Độ dài khoảng tin cậy}]
Độ dài $L = \hat{\Theta}_U - \hat{\Theta}_L$. Cùng mẫu, độ tin cậy cao hơn $\Rightarrow$ khoảng rộng hơn. Cỡ mẫu lớn hơn $\Rightarrow$ khoảng hẹp hơn (vì $\text{SE}(\bar{X}) = \sigma/\sqrt{n}$).
\end{ynghia}

\subsubsection{KTC cho kỳ vọng - phương sai đã biết}

\begin{congthuc}[title={KTC đối xứng cho $\mu$ khi $\sigma$ biết}]
  $X \sim \mathcal{N}(\mu,\sigma^2)$, $\sigma$ biết, mẫu cỡ $n$:
  \[
    \bar{X} - z_{\frac{\alpha}{2}}\,\dfrac{\sigma}{\sqrt{n}} \le \mu \le \bar{X} + z_{\frac{\alpha}{2}}\,\dfrac{\sigma}{\sqrt{n}}.
  \]
  Với $1-\alpha = 95\%$: $z_{0{,}025} = 1{,}96$. Với $99\%$: $z_{0{,}005} = 2{,}58$.
\end{congthuc}

\begin{vidu}[title={Trọng lượng sản phẩm - KTC cho $\mu$}]
Trọng lượng $X \sim \mathcal{N}(\mu, \sigma^2)$, $\sigma = 6$ gam. Cân $n=36$ sản phẩm, $\bar{x} = 166{,}22$ gam.

\medskip
\textbf{(a)} Độ tin cậy 95\% ($z_{0{,}025} = 1{,}96$):
\[
  166{,}22 - 1{,}96 \cdot \dfrac{6}{\sqrt{36}} \le \mu \le 166{,}22 + 1{,}96 \cdot \dfrac{6}{\sqrt{36}}
\]
\[
  \boxed{164{,}26 \le \mu \le 168{,}18 \text{ (gam)}}.
\]

\medskip
\textbf{(b)} Độ tin cậy 99\% ($z_{0{,}005} = 2{,}58$):
\[
  \boxed{163{,}64 \le \mu \le 168{,}80 \text{ (gam)}}.
\]
Cùng mẫu, độ tin cậy lớn hơn $\Rightarrow$ khoảng rộng hơn.
\end{vidu}

\subsubsection{KTC cho kỳ vọng - phương sai chưa biết}

\begin{congthuc}[title={KTC cho $\mu$ khi $\sigma$ chưa biết}]
  $X \sim \mathcal{N}(\mu,\sigma^2)$, $\sigma$ chưa biết, $n$ nhỏ:
  \[
    \bar{X} - t_{\frac{\alpha}{2}}(n-1)\,\dfrac{S}{\sqrt{n}} \le \mu \le \bar{X} + t_{\frac{\alpha}{2}}(n-1)\,\dfrac{S}{\sqrt{n}}.
  \]
  Thống kê $T = \dfrac{\bar{X}-\mu}{S/\sqrt{n}} \sim t(n-1)$.
\end{congthuc}

\begin{vidu}[title={KTC cho trọng lượng - $\sigma$ chưa biết}]
$n=25$, $\bar{x}=50{,}2$, $s=4{,}8$. KTC 95\%: $t_{0{,}025;24} = 2{,}064$.
\[
  50{,}2 \pm 2{,}064 \cdot \dfrac{4{,}8}{\sqrt{25}} \;\Rightarrow\; [48{,}22;\; 52{,}18].
\]
\end{vidu}

\begin{luuy}[title={Khi nào dùng $z$ và $t$?}]
\begin{itemize}
  \item $\sigma$ biết (hiếm): dùng $z$.
  \item $\sigma$ chưa biết, $X$ chuẩn, $n$ nhỏ: dùng $t(n-1)$.
  \item $n \ge 30$, không chắc chuẩn: CLT $\Rightarrow$ dùng $z$ với $S$ thay $\sigma$.
\end{itemize}
\end{luuy}

\subsubsection{KTC cho phương sai}

\begin{congthuc}[title={KTC cho $\sigma^2$}]
  $X \sim \mathcal{N}(\mu,\sigma^2)$, mẫu cỡ $n$:
  \[
    \dfrac{(n-1)S^2}{\chi^2_{\frac{\alpha}{2}}(n-1)} \le \sigma^2 \le \dfrac{(n-1)S^2}{\chi^2_{1-\frac{\alpha}{2}}(n-1)}.
  \]
  KTC cho $\sigma$: lấy căn bậc hai hai cận (không đối xứng nếu lấy trên $\sigma^2$).
\end{congthuc}

\begin{vidu}[title={KTC cho phương sai}]
$n=16$, $s^2 = 25$. KTC 95\% cho $\sigma^2$:
\[
  \dfrac{15 \cdot 25}{\chi^2_{0{,}025;15}} \le \sigma^2 \le \dfrac{15 \cdot 25}{\chi^2_{0{,}975;15}}
  \;\Rightarrow\; [13{,}33;\; 54{,}87].
\]
\end{vidu}

\subsubsection{KTC cho tỷ lệ}

\begin{congthuc}[title={KTC cho tỷ lệ $p$ - mẫu lớn}]
  $X_i \sim \text{Bernoulli}(p)$, độc lập cùng phân phối; $n$ lớn ($np \ge 5$, $n(1-p) \ge 5$):
  \[
    \hat{p} - z_{\frac{\alpha}{2}}\sqrt{\dfrac{\hat{p}(1-\hat{p})}{n}} \le p \le \hat{p} + z_{\frac{\alpha}{2}}\sqrt{\dfrac{\hat{p}(1-\hat{p})}{n}},
  \]
  với $\hat{p} = \dfrac{x}{n}$, $x$ = số phần tử có dấu hiệu A.
\end{congthuc}

\begin{vidu}[title={KTC cho tỷ lệ phế phẩm 4\%}]
Nhà máy: $n=500$, $x=20$, $\hat{p} = 0{,}04$. KTC 95\%:
\[
  0{,}04 \pm 1{,}96\sqrt{\dfrac{0{,}04 \cdot 0{,}96}{500}}
  = 0{,}04 \pm 0{,}0172 \;\Rightarrow\; [0{,}0228;\; 0{,}0572].
\]
Tức tỷ lệ phế phẩm thật có thể từ 2{,}28\% đến 5{,}72\% (95\% lần lấy mẫu tương tự).
\end{vidu}

\begin{vidu}[title={KTC cho tỷ lệ}]
Khảo sát 400 người, 120 ủng hộ chính sách mới: $\hat{p} = 0{,}3$. KTC 95\%:
\[
  0{,}3 \pm 1{,}96\sqrt{\dfrac{0{,}3 \cdot 0{,}7}{400}} \;\Rightarrow\; [0{,}255;\; 0{,}345].
\]
\end{vidu}

\subsubsection{Sai số ước lượng và cỡ mẫu}

\begin{dinhnghia}[title={Sai số ước lượng}]
  $\hat{\Theta}$ ước lượng $\theta$. \textbf{Sai số} $\varepsilon$ với độ tin cậy $1-\alpha$ nếu $P(|\hat{\Theta}-\theta| \le \varepsilon) \ge 1-\alpha$.
  Với $\mu$, $\sigma$ biết: $\varepsilon = \dfrac{z_{\frac{\alpha}{2}}\,\sigma}{\sqrt{n}}$.
\end{dinhnghia}

\begin{vidu}[title={Sai số}]
Ví dụ 25(a): sai số $\varepsilon = \dfrac{1{,}96 \cdot 6}{\sqrt{36}} = 1{,}96$ gam với độ tin cậy 95\%.
\end{vidu}

\begin{congthuc}[title={Cỡ mẫu cần thiết}]
  \textbf{Cho $\mu$} ($\sigma$ biết), sai số $\le \varepsilon$:
  \[
    n \ge \left(\dfrac{z_{\frac{\alpha}{2}}\,\sigma}{\varepsilon}\right)^2.
  \]
  \textbf{Cho $p$}, sai số $\le \varepsilon$:
  \[
    n \ge \dfrac{z_{\frac{\alpha}{2}}^2\,p(1-p)}{\varepsilon^2}.
  \]
  Nếu chưa biết $p$, dùng $p = 0{,}5$ (conservative, cho $n$ lớn nhất).
\end{congthuc}

\begin{vidu}[title={Cỡ mẫu cho}]
Muốn ước lượng $\mu$ với sai số $\varepsilon = 1$ gam, độ tin cậy 95\%, $\sigma = 6$:
\[
  n \ge \left(\dfrac{1{,}96 \cdot 6}{1}\right)^2 = 138{,}3 \;\Rightarrow\; \text{chọn } n = 139.
\]
\end{vidu}

\begin{vidu}[title={Cỡ mẫu cho tỷ lệ}]
Ước lượng $p$ với sai số $\varepsilon = 0{,}02$, độ tin cậy 95\%, chưa biết $p$:
\[
  n \ge \dfrac{1{,}96^2 \cdot 0{,}25}{0{,}02^2} = 2401.
\]
\end{vidu}

\begin{vidu}[title={Cận trên một phía}]
Ví dụ 25, độ tin cậy 99\%, cận trên phải cho $\mu$:
\[
  \mu \le \bar{x} + z_{\alpha}\,\dfrac{\sigma}{\sqrt{n}} = 166{,}22 + 2{,}326 \cdot 1 = 168{,}55 \text{ gam}.
\]
\end{vidu}

\begin{phuongphap}[title={Quy trình lập KTC}]
\begin{enumerate}
  \item Xác định tham số $\theta$, độ tin cậy $1-\alpha$.
  \item Chọn công thức phù hợp ($\mu$, $\sigma^2$, $p$; $\sigma$ biết/chưa biết).
  \item Tính thống kê mẫu ($\bar{x}$, $s$, $\hat{p}$).
  \item Tra $z$, $t$, hoặc $\chi^2$; thay số; diễn giải bằng lời.
\end{enumerate}
\end{phuongphap}

\begin{tongket}[title={Tổng kết}]
\begin{enumerate}
  \item KTC mức $1-\alpha$: $P(\hat{\Theta}_L \le \theta \le \hat{\Theta}_U) \ge 1-\alpha$.
  \item $\mu$, $\sigma$ biết: $\bar{X} \pm \dfrac{z_{\frac{\alpha}{2}}\,\sigma}{\sqrt{n}}$.
  \item $\mu$, $\sigma$ chưa biết: $\bar{X} \pm \dfrac{t_{\frac{\alpha}{2}}(n-1)\,S}{\sqrt{n}}$.
  \item $\sigma^2$: dùng $\chi^2(n-1)$; $p$: $\hat{p} \pm z_{\frac{\alpha}{2}}\sqrt{\dfrac{\hat{p}(1-\hat{p})}{n}}$.
  \item Cỡ mẫu: $n \ge (\dfrac{z_{\frac{\alpha}{2}}\sigma}{\varepsilon})^2$.
  \item Làm lại Ví dụ 25: $[164{,}26;\;168{,}18]$.
\end{enumerate}
\end{tongket}

\begin{tongket}[title={Tổng kết}]
\begin{enumerate}
  \item Mẫu độc lập cùng phân phối, thống kê mẫu, phân phối mẫu, CLT, $\chi^2$, $t$, $F$.
  \item Ước lượng điểm; không chệch, vững, MSE; MME, MLE.
  \item Khoảng tin cậy cho $\mu$, $\sigma^2$, $p$; cỡ mẫu.
\end{enumerate}
\end{tongket}

\newpage
\section{Kiểm định giả thuyết thống kê}

\begin{dongco}[title={Động cơ}]
Chương trước trả lời ``$\theta$ khoảng bao nhiêu?''. Chương này trả lời ``$\theta$ có \textbf{bằng} giá trị quy định không?'' - ví dụ máy còn sản xuất đúng 100 gam hay đã lệch?
\end{dongco}

\subsection{Giả thuyết thống kê và quy tắc kiểm định}

\subsubsection{Giả thuyết thống kê}

\begin{dinhnghia}[title={Giả thuyết thống kê}]
  Mọi khẳng định về tham số, dạng phân phối, hoặc tính độc lập của biến ngẫu nhiên đều gọi là \textbf{giả thuyết thống kê}. Tập trung vào giả thuyết về tham số ($\mu$, $\sigma^2$, $p$, $\mu_1-\mu_2$, \ldots).
\end{dinhnghia}

\begin{dinhnghia}[title={Giả thuyết không và đối thuyết}]
  \textbf{Giả thuyết không} $H_0$: giả thuyết cần kiểm tra (thường ``không có thay đổi / không khác biệt'').
  
  \textbf{Giả thuyết đối} $H_1$: khẳng định ngược lại (có thay đổi / có khác biệt).
\end{dinhnghia}

\begin{vidu}[title={Cặp giả thuyết về $\mu$}]
\begin{itemize}
  \item \textbf{Hai phía:} $H_0: \mu = \mu_0$ đối với $H_1: \mu \ne \mu_0$.
  \item \textbf{Một phía phải:} $H_0: \mu = \mu_0$ đối với $H_1: \mu > \mu_0$.
  \item \textbf{Một phía trái:} $H_0: \mu = \mu_0$ đối với $H_1: \mu < \mu_0$.
\end{itemize}
\end{vidu}

\begin{vidu}[title={Giả thuyết về tỷ lệ phế phẩm}]
Nhà máy quy định tỷ lệ phế phẩm $\le 3\%$. Kiểm tra có vượt không:
\[
  H_0: p = 0{,}03 \quad\text{vs}\quad H_1: p > 0{,}03.
\]
\end{vidu}

\subsubsection{Tiêu chuẩn kiểm định và miền bác bỏ}

\begin{dinhnghia}[title={Tiêu chuẩn kiểm định}]
  \textbf{Tiêu chuẩn kiểm định} $T_0 = T(X_1,\ldots,X_n)$ là thống kê dùng để quyết định bác bỏ $H_0$ hay không. Giá trị quan sát $t_0$ tính từ mẫu cụ thể.
\end{dinhnghia}

\begin{dinhnghia}[title={Miền bác bỏ và mức ý nghĩa}]
  \textbf{Miền bác bỏ} $W_\alpha$: tập giá trị của $T_0$ khiến ta \textbf{bác bỏ} $H_0$.
  
  \textbf{Mức ý nghĩa} $\alpha$: xác suất \textbf{sai lầm loại I} - bác bỏ $H_0$ khi $H_0$ đúng.
  
  Quy tắc: nếu $t_0 \in W_\alpha$ thì bác bỏ $H_0$; ngược lại \textbf{chưa đủ cơ sở} bác bỏ $H_0$.
\end{dinhnghia}

\begin{congthuc}[title={Miền bác bỏ theo dạng $H_1$}]
  \begin{center}
  \begin{tabular}{lll}
    \toprule
    $H_0$ & $H_1$ & $W_\alpha$ \\
    \midrule
    $\mu = \mu_0$ & $\mu \ne \mu_0$ & $(-\infty, -z_{\frac{\alpha}{2}}) \cup (z_{\frac{\alpha}{2}}, +\infty)$ \\
    $\mu = \mu_0$ & $\mu > \mu_0$ & $(z_\alpha, +\infty)$ \\
    $\mu = \mu_0$ & $\mu < \mu_0$ & $(-\infty, -z_\alpha)$ \\
    \bottomrule
  \end{tabular}
  \end{center}
\end{congthuc}

\begin{ynghia}[title={Giá trị $p$}]
\textbf{Giá trị $p$} (p-value trong tài liệu tiếng Anh) là xác suất, giả sử $H_0$ đúng, thu được thống kê ``cực đoan hơn hoặc bằng'' giá trị quan sát. Nếu giá trị $p < \alpha$ thì bác bỏ $H_0$. Giá trị $p$ càng nhỏ, bằng chứng chống $H_0$ càng mạnh.
\end{ynghia}

\subsubsection{Sai lầm loại I và loại II}

\begin{dinhnghia}[title={Sai lầm kiểm định}]
\begin{itemize}
  \item \textbf{Loại I} ($\alpha$): bác bỏ $H_0$ khi $H_0$ \textbf{đúng} (``báo động giả'').
  \item \textbf{Loại II} ($\beta$): \textbf{không} bác bỏ $H_0$ khi $H_0$ \textbf{sai} (``bỏ sót'').
  \item \textbf{Lực kiểm định}: $1 - \beta$ = xác suất bác bỏ $H_0$ khi $H_0$ sai.
\end{itemize}
\end{dinhnghia}

\begin{congthuc}[title={Bảng quyết định}]
  \begin{center}
  \begin{tabular}{lcc}
    \toprule
    Quyết định & $H_0$ đúng & $H_0$ sai \\
    \midrule
    Bác bỏ $H_0$ & Sai lầm I ($\alpha$) & Đúng ($1-\beta$) \\
    Không bác bỏ $H_0$ & Đúng ($1-\alpha$) & Sai lầm II ($\beta$) \\
    \bottomrule
  \end{tabular}
  \end{center}
\end{congthuc}

\begin{luuy}[title={Đánh đổi $\alpha$ và $\beta$}]
Giảm $\alpha$ (khắt khe hơn với $H_0$) thường làm $\beta$ tăng (dễ bỏ sót). Tăng cỡ mẫu $n$ giúp \textbf{cả} giảm $\beta$ lẫn thu hẹp KTC - ``mua'' độ chính xác bằng chi phí lấy mẫu.
\end{luuy}

\begin{vidu}[title={Ví dụ tư pháp - $\alpha$ và $\beta$}]
Xét nghiệm tội phạm: $H_0$ = vô tội, $H_1$ = có tội. Sai lầm I = kết tội người vô tội; sai lầm II = tha người có tội. Xã hội thường chọn $\alpha$ rất nhỏ (``vô tội cho đến khi chứng minh ngược lại'').
\end{vidu}

\subsubsection{Các bước kiểm định giả thuyết}

\begin{phuongphap}[title={Quy trình 5 bước}]
\begin{enumerate}
  \item \textbf{Lập} cặp giả thuyết $H_0$, $H_1$ (theo yêu cầu bài toán).
  \item \textbf{Chọn} tiêu chuẩn kiểm định $T_0$ và xác định phân phối mẫu khi $H_0$ đúng.
  \item \textbf{Xây dựng} miền bác bỏ $W_\alpha$ với mức ý nghĩa $\alpha$ cho trước.
  \item \textbf{Tính} giá trị quan sát $t_0$ từ mẫu.
  \item \textbf{Kết luận:} $t_0 \in W_\alpha$ $\Rightarrow$ bác bỏ $H_0$; ngược lại chưa đủ cơ sở bác bỏ $H_0$.
\end{enumerate}
\end{phuongphap}

\begin{luuy}[title={Diễn đạt kết luận đúng}]
Không nói ``chấp nhận $H_0$'' mà nói ``\textbf{chưa có cơ sở để bác bỏ} $H_0$''. Không bác bỏ $H_0$ $\not=$ chứng minh $H_0$ đúng - có thể do $n$ nhỏ, $\beta$ lớn.
\end{luuy}

\begin{vidu}[title={Kiểm định hai phía - minh họa}]
$H_0: \mu = 10$ đối với $H_1: \mu \ne 10$, $\alpha = 5\%$, $Z_0 = \dfrac{\bar{X}-10)}{\sigma/\sqrt{n}}$.
$W_{0{,}05} = (-\infty, -1{,}96) \cup (1{,}96, +\infty)$.
Nếu $z_0 = 1{,}5 \notin W_{0{,}05}$: chưa đủ cơ sở bác bỏ $H_0$.
\end{vidu}

\begin{tongket}[title={Tổng kết}]
\begin{enumerate}
  \item $H_0$ = giả thuyết cần kiểm tra; $H_1$ = đối thuyết.
  \item Tiêu chuẩn kiểm định + miền bác bỏ $W_\alpha$.
  \item Sai lầm I ($\alpha$), II ($\beta$); lực $1-\beta$.
  \item 5 bước giải bài toán kiểm định.
\end{enumerate}
\end{tongket}
\subsection{Kiểm định giả thuyết về tham số một tổng thể}

\subsubsection{Kiểm định kỳ vọng - $\sigma$ đã biết}

\begin{dongco}[title={Động cơ - Máy sản xuất có lệch không?}]
Máy quy định $\mu_0 = 100$ gam, $\sigma = 2$ gam. Cân $n=100$, $\bar{x} = 100{,}4$. Đây có phải dao động ngẫu nhiên hay dấu hiệu máy đã tăng trọng lượng?
\end{dongco}

\begin{congthuc}[title={Kiểm định $\mu$ khi $\sigma$ biết}]
  $X \sim \mathcal{N}(\mu,\sigma^2)$, $\sigma$ biết:
  \[
    Z_0 = \dfrac{\bar{X} - \mu_0}{\sigma/\sqrt{n}}.
  \]
  Nếu $H_0: \mu = \mu_0$ đúng thì $Z_0 \sim \mathcal{N}(0,1)$.
\end{congthuc}

\begin{vidu}[title={Kiểm định trọng lượng máy - kiểm định $z$}]
$X \sim \mathcal{N}(\mu, 2^2)$. Nghi ngờ trọng lượng \textbf{tăng}: $H_0: \mu = 100$ đối với $H_1: \mu > 100$, $\alpha = 5\%$.

\medskip
Tiêu chuẩn: $Z_0 = \dfrac{\bar{X} - 100}{2/\sqrt{n}}$. Khi $H_0$ đúng: $Z_0 \sim \mathcal{N}(0,1)$.

Miền bác bỏ (một phía phải): $W_{0{,}05} = (z_{0{,}05}; +\infty) = (1{,}645; +\infty)$.

\medskip
$n=100$, $\bar{x} = 100{,}4$:
\[
  z_0 = \dfrac{100{,}4 - 100}{2/\sqrt{100}} = \dfrac{0{,}4}{0{,}2} = 2.
\]
Vì $z_0 = 2 \in W_{0{,}05}$ nên \textbf{bác bỏ} $H_0$: có cơ sở cho rằng trọng lượng trung bình đã tăng (mức ý nghĩa 5\%).

\medskip
\textbf{Giá trị $p$} (một phía phải): $P(Z > 2) = 1 - \Phi(2) \approx 0{,}0228 < 0{,}05$.
\end{vidu}

\begin{congthuc}[title={Miền bác bỏ cho kiểm định $\mu$, $\sigma$ biết}]
  \begin{center}
  \begin{tabular}{lll}
    \toprule
    $H_1$ & $W_\alpha$ & Ví dụ $\alpha=5\%$ \\
    \midrule
    $\mu \ne \mu_0$ & $|Z_0| > z_{\frac{\alpha}{2}}$ & $|Z_0| > 1{,}96$ \\
    $\mu > \mu_0$ & $Z_0 > z_\alpha$ & $Z_0 > 1{,}645$ \\
    $\mu < \mu_0$ & $Z_0 < -z_\alpha$ & $Z_0 < -1{,}645$ \\
    \bottomrule
  \end{tabular}
  \end{center}
\end{congthuc}

\subsubsection{Kiểm định kỳ vọng - mẫu kích thước lớn}

\begin{phuongphap}[title={Mẫu lớn, $\sigma$ chưa biết}]
  Khi $n \ge 30$, không chắc phân phối chuẩn, thay $\sigma$ bằng $S$:
  \[
    Z_0 = \dfrac{\bar{X} - \mu_0}{S/\sqrt{n}} \approx \mathcal{N}(0,1) \quad \text{dưới } H_0.
  \]
  Miền bác bỏ giống trường hợp $Z$ chuẩn.
\end{phuongphap}

\begin{vidu}[title={Kiểm định mẫu lớn}]
$n=64$, $\bar{x}=52$, $s=8$, $H_0:\mu=50$ đối với $H_1:\mu \ne 50$, $\alpha=5\%$:
\[
  z_0 = \dfrac{52-50}{8/\sqrt{64}} = \dfrac{2}{1} = 2 > 1{,}96 \;\Rightarrow\; \text{bác bỏ } H_0.
\]
\end{vidu}

\subsubsection{Kiểm định kỳ vọng - $\sigma$ chưa biết, $n$ nhỏ}

\begin{congthuc}[title={Kiểm định $t$}]
  $X \sim \mathcal{N}(\mu,\sigma^2)$, $\sigma$ chưa biết, $n < 30$:
  \[
    T_0 = \dfrac{\bar{X} - \mu_0}{S/\sqrt{n}} \sim t(n-1) \quad \text{dưới } H_0.
  \]
  Miền bác bỏ hai phía: $|T_0| > t_{\frac{\alpha}{2}}(n-1)$.
\end{congthuc}

\begin{vidu}[title={Kiểm định năng suất hạt giống - kiểm định $t$}]
Công ty tuyên bố năng suất trung bình $\mu_0 = 21{,}5$ tạ/ha. Gieo thử $n=16$ vườn:
\[
  19{,}2;\; 18{,}7;\; 22{,}4;\; 20{,}3;\; 16{,}8;\; 25{,}1;\; 17{,}0;\; 15{,}8;\; 21{,}0;\; 18{,}6;\; 23{,}7;\; 24{,}1;\; 23{,}4;\; 19{,}8;\; 21{,}7;\; 18{,}9.
\]
$X \sim \mathcal{N}(\mu,\sigma^2)$, $\sigma$ chưa biết.

\medskip
$H_0: \mu = 21{,}5$ đối với $H_1: \mu \ne 21{,}5$, $\alpha = 5\%$.

Tiêu chuẩn: $T_0 = \dfrac{\bar{X} - 21{,}5}{S/\sqrt{16}}$. Dưới $H_0$: $T_0 \sim t(15)$.

Miền bác bỏ: $W_{0{,}05} = (-\infty, -2{,}131) \cup (2{,}131, +\infty)$ (vì $t_{0{,}025;15} = 2{,}131$).

\medskip
Từ dữ liệu: $\bar{x} = 20{,}406$, $s = 3{,}038$.
\[
  t_0 = \dfrac{20{,}406 - 21{,}5}{3{,}038/\sqrt{16}} = \dfrac{-1{,}094}{0{,}7595} \approx -1{,}4404.
\]
Vì $t_0 = -1{,}4404 \notin W_{0{,}05}$ nên \textbf{chưa có cơ sở bác bỏ} $H_0$: có thể chấp nhận lời quảng cáo 21{,}5 tạ/ha (mức ý nghĩa 5\%).
\end{vidu}

\begin{luuy}[title={Chọn $z$ hay $t$}]
\begin{itemize}
  \item $\sigma$ biết: luôn dùng $Z_0$.
  \item $\sigma$ chưa biết, $n < 30$, $X$ chuẩn: dùng $T_0 \sim t(n-1)$.
  \item $\sigma$ chưa biết, $n \ge 30$: dùng $Z_0$ với $S$ (xấp xỉ chuẩn).
\end{itemize}
\end{luuy}

\subsubsection{Kiểm định phương sai}

\begin{congthuc}[title={Kiểm định $\chi^2$ cho $\sigma^2$}]
  $H_0: \sigma^2 = \sigma_0^2$ đối với $H_1: \sigma^2 \ne \sigma_0^2$ (hoặc một phía):
  \[
    \chi^2_0 = \dfrac{(n-1)S^2}{\sigma_0^2} \sim \chi^2(n-1) \quad \text{dưới } H_0.
  \]
  Hai phía: bác bỏ nếu $\chi^2_0 < \chi^2_{1-\frac{\alpha}{2}}(n-1)$ hoặc $\chi^2_0 > \chi^2_{\frac{\alpha}{2}}(n-1)$.
\end{congthuc}

\begin{vidu}[title={Kiểm định phương sai máy cân}]
Máy cân có $\sigma_0 = 0{,}5$ gam. Kiểm tra 20 sản phẩm, $s = 0{,}62$. $H_0: \sigma = 0{,}5$ đối với $H_1: \sigma \ne 0{,}5$, $\alpha=5\%$:
\[
  \chi^2_0 = \dfrac{19 \cdot 0{,}62^2}{0{,}5^2} = \dfrac{19 \cdot 0{,}3844}{0{,}25} \approx 29{,}2.
\]
$\chi^2_{0{,}025;19} = 32{,}852$, $\chi^2_{0{,}975;19} = 8{,}907$. Vì $8{,}907 < 29{,}2 < 32{,}852$, \textbf{chưa bác bỏ} $H_0$.
\end{vidu}

\subsubsection{Kiểm định tỷ lệ}

\begin{congthuc}[title={Kiểm định tỷ lệ - mẫu lớn}]
  $H_0: p = p_0$ đối với $H_1: p \ne p_0$ (hoặc một phía), $n$ lớn:
  \[
    Z_0 = \dfrac{\hat{P} - p_0}{\sqrt{\dfrac{p_0(1-p_0)}{n}}} \approx \mathcal{N}(0,1) \quad \text{dưới } H_0.
  \]
\end{congthuc}

\begin{vidu}[title={Kiểm định tỷ lệ phế phẩm 4\%}]
Nhà máy cam kết $p_0 = 3\%$. Mẫu $n=500$, $x=20$, $\hat{p}=4\%$. $H_0: p=0{,}03$ đối với $H_1: p>0{,}03$, $\alpha=5\%$:
\[
  z_0 = \dfrac{0{,}04 - 0{,}03}{\sqrt{\dfrac{0{,}03 \cdot 0{,}97}{500}}} = \dfrac{0{,}01}{0{,}00783} \approx 1{,}28.
\]
$z_{0{,}05} = 1{,}645$. Vì $1{,}28 < 1{,}645$, \textbf{chưa bác bỏ} $H_0$ (chưa đủ bằng chứng tỷ lệ vượt 3\%).
\end{vidu}

\begin{vidu}[title={Kiểm định tỷ lệ}]
Thử nghiệm thuốc: 200 bệnh nhân, 160 khỏi. $H_0: p=0{,}75$ đối với $H_1: p>0{,}75$:
\[
  z_0 = \dfrac{0{,}80 - 0{,}75}{\sqrt{\dfrac{0{,}75 \cdot 0{,}25}{200}}} = \dfrac{0{,}05}{0{,}0306} \approx 1{,}63.
\]
$z_{0{,}05} = 1{,}645$: sát ngưỡng; giá trị $p \approx 0{,}051 > 0{,}05$, chưa bác bỏ $H_0$.
\end{vidu}

\begin{congthuc}[title={Tóm tắt kiểm định một mẫu}]
  \begin{center}
  \begin{tabular}{llll}
    \toprule
    Tham số & $H_0$ & Tiêu chuẩn & Phân phối ($H_0$ đúng) \\
    \midrule
    $\mu$ ($\sigma$ biết) & $\mu=\mu_0$ & $Z_0 = \dfrac{\bar{X}-\mu_0}{\sigma/\sqrt{n}}$ & $\mathcal{N}(0,1)$ \\
    $\mu$ ($\sigma$ ?) & $\mu=\mu_0$ & $T_0 = \dfrac{\bar{X}-\mu_0}{S/\sqrt{n}}$ & $t(n-1)$ hoặc $\approx\mathcal{N}(0,1)$ \\
    $\sigma^2$ & $\sigma^2=\sigma_0^2$ & $\chi^2_0 = \dfrac{(n-1)S^2}{\sigma_0^2}$ & $\chi^2(n-1)$ \\
    $p$ & $p=p_0$ & $Z_0 = \dfrac{\hat{P}-p_0}{\sqrt{\dfrac{p_0(1-p_0)}{n}}}$ & $\approx\mathcal{N}(0,1)$ \\
    \bottomrule
  \end{tabular}
  \end{center}
\end{congthuc}

\begin{tongket}[title={Tổng kết}]
\begin{enumerate}
  \item Kiểm định $z$ ($\sigma$ biết): Ví dụ máy 100 gam, $\bar{x}=100{,}4$, $z_0=2$, bác bỏ $H_0$.
  \item Kiểm định $t$ ($\sigma$ chưa biết, $n$ nhỏ): Ví dụ hạt giống, $t_0 \approx -1{,}44$, chưa bác bỏ $H_0$.
  \item Kiểm định $\chi^2$ cho $\sigma^2$; kiểm định $z$ cho $p$.
  \item Thuộc 5 bước kiểm định và miền bác bỏ theo dạng $H_1$.
\end{enumerate}
\end{tongket}
\subsection{So sánh hai tổng thể}

\subsubsection{So sánh hai kỳ vọng - phương sai đã biết}

\begin{dinhnghia}[title={Bài toán}]
  Hai tổng thể, $X_i \sim \mathcal{N}(\mu_i, \sigma_i^2)$, mẫu độc lập cỡ $n_1$, $n_2$. Kiểm định:
  \[
    H_0: \mu_1 - \mu_2 = \Delta_0 \quad \text{vs} \quad H_1: \mu_1 - \mu_2 \ne \Delta_0
  \]
  ($\Delta_0 = 0$ nếu so sánh ``có khác nhau không'').
\end{dinhnghia}

\begin{congthuc}[title={Tiêu chuẩn kiểm định, $\sigma_1, \sigma_2$ biết}]
  \[
    Z_0 = \dfrac{(\bar{X}_1 - \bar{X}_2) - \Delta_0}{\sqrt{\dfrac{\sigma_1^2}{n_1} + \dfrac{\sigma_2^2}{n_2}}} \sim \mathcal{N}(0,1) \quad \text{dưới } H_0.
  \]
\end{congthuc}

\begin{vidu}[title={So sánh hai dây chuyền sản xuất}]
Dây chuyền I: $n_1=16$, $\bar{x}_1=100$, $\sigma_1=2$. Dây chuyền II: $n_2=20$, $\bar{x}_2=102$, $\sigma_2=2{,}5$.
$H_0: \mu_1=\mu_2$ đối với $H_1: \mu_1 \ne \mu_2$, $\alpha=5\%$:
\[
  z_0 = \dfrac{(100-102)-0}{\sqrt{\dfrac{4}{16} + \dfrac{6{,}25}{20}}} = \dfrac{-2}{\sqrt{0{,}5625}} = \dfrac{-2}{0{,}75} \approx -2{,}667.
\]
$|z_0| > 1{,}96$ $\Rightarrow$ \textbf{bác bỏ} $H_0$: hai dây chuyền có trung bình khác nhau.
\end{vidu}

\subsubsection{So sánh hai kỳ vọng - mẫu lớn}

\begin{phuongphap}[title={Hai mẫu lớn, $\sigma$ chưa biết}]
  Thay $\sigma_i$ bằng $S_i$, $n_1, n_2 \ge 30$:
  \[
    Z_0 = \dfrac{(\bar{X}_1 - \bar{X}_2) - \Delta_0}{\sqrt{\dfrac{S_1^2}{n_1} + \dfrac{S_2^2}{n_2}}} \approx \mathcal{N}(0,1).
  \]
\end{phuongphap}

\subsubsection{So sánh hai kỳ vọng - phương sai chưa biết}

\begin{congthuc}[title={Hai mẫu độc lập, phương sai chưa biết}]
  \textbf{Trường hợp $\sigma_1^2 = \sigma_2^2$:}
  \[
    S_p^2 = \dfrac{(n_1-1)S_1^2 + (n_2-1)S_2^2}{n_1+n_2-2}, \quad
    T_0 = \dfrac{(\bar{X}_1-\bar{X}_2)-\Delta_0}{S_p\sqrt{\dfrac{1}{n_1}+\dfrac{1}{n_2}}} \sim t(n_1+n_2-2).
  \]
  \textbf{Trường hợp $\sigma_1^2 \ne \sigma_2^2$:} dùng phân phối xấp xỉ.
\end{congthuc}

\begin{vidu}[title={So sánh năng suất hai chất xúc tác}]
Chất xúc tác I: $n_1=10$, $\bar{x}_1=85$, $s_1=5$. Chất II: $n_2=12$, $\bar{x}_2=80$, $s_2=6$.
$H_0: \mu_1=\mu_2$ đối với $H_1: \mu_1 \ne \mu_2$, giả sử $\sigma_1^2=\sigma_2^2$, $\alpha=5\%$:
\[
  s_p^2 = \dfrac{9\cdot 25 + 11\cdot 36}{20} = \dfrac{621}{20} = 31{,}05, \quad s_p \approx 5{,}573.
\]
\[
  t_0 = \dfrac{85-80}{5{,}573\sqrt{\dfrac{1}{10}+\dfrac{1}{12}}} = \dfrac{5}{5{,}573 \cdot 0{,}429} \approx 2{,}09.
\]
$t_{0{,}025;20} = 2{,}086$. Vì $|t_0| \approx 2{,}09 > 2{,}086$ nên \textbf{bác bỏ} $H_0$: hai chất xúc tác cho năng suất trung bình khác nhau.
\end{vidu}

\subsubsection{So sánh cặp đôi}

\begin{dongco}[title={Động cơ - Cùng đối tượng, hai điều kiện}]
Đo cùng 10 bệnh nhân trước và sau uống thuốc; cùng 12 thửa ruộng trước và sau bón phân. Hai mẫu \textbf{phụ thuộc} - không dùng công thức hai mẫu độc lập.
\end{dongco}

\begin{phuongphap}[title={Kiểm định t cặp đôi}]
  Tính hiệu $D_i = X_{1i} - X_{2i}$, $i=1,\ldots,n$. Coi $D_i$ độc lập, cùng phân phối $\mathcal{N}(\mu_D, \sigma_D^2)$.
  \[
    H_0: \mu_D = \Delta_0 \quad \text{vs} \quad H_1: \mu_D \ne \Delta_0.
  \]
  Tiêu chuẩn:
  \[
    T_0 = \dfrac{\bar{D} - \Delta_0}{S_D/\sqrt{n}} \sim t(n-1) \quad \text{dưới } H_0.
  \]
\end{phuongphap}

\begin{vidu}[title={Kiểm định cặp - chế độ chăm bón}]
12 thửa ruộng, năng suất (tạ/ha) trước/sau bón phân:
\begin{center}
\begin{tabular}{cccccccccccc}
\toprule
Thửa & 1 & 2 & 3 & 4 & 5 & 6 & 7 & 8 & 9 & 10 & 11 & 12 \\
\midrule
Trước & 42 & 38 & 45 & 40 & 44 & 39 & 41 & 43 & 37 & 46 & 40 & 42 \\
Sau & 48 & 44 & 50 & 46 & 49 & 45 & 47 & 48 & 43 & 51 & 45 & 47 \\
\bottomrule
\end{tabular}
\end{center}
Hiệu $d_i$: 6, 6, 5, 6, 5, 6, 6, 5, 6, 5, 5, 5. $\bar{d} = 5{,}583$, $s_d \approx 0{,}515$.

$H_0: \mu_D = 0$ đối với $H_1: \mu_D > 0$, $\alpha=1\%$:
\[
  t_0 = \dfrac{5{,}583}{0{,}515/\sqrt{12}} \approx 37{,}5 \gg t_{0{,}01;11} = 2{,}718.
\]
\textbf{Bác bỏ} $H_0$: bón phân làm tăng năng suất có ý nghĩa thống kê.
\end{vidu}

\begin{vidu}[title={Kiểm định cặp - huyết áp trước/sau thuốc}]
10 bệnh nhân, huyết áp tâm thu (mmHg) trước/sau: hiệu trung bình $\bar{d} = -8$, $s_d = 4$, $n=10$.
$H_0: \mu_D = 0$ đối với $H_1: \mu_D < 0$:
\[
  t_0 = \dfrac{-8}{4/\sqrt{10}} = -6{,}32.
\]
$|t_0| \gg t_{0{,}025;9} = 2{,}262$ $\Rightarrow$ thuốc giảm huyết áp có ý nghĩa.
\end{vidu}

\subsubsection{So sánh hai phương sai}

\begin{congthuc}[title={Kiểm định F}]
  $H_0: \sigma_1^2 = \sigma_2^2$ đối với $H_1: \sigma_1^2 \ne \sigma_2^2$:
  \[
    F_0 = \dfrac{S_1^2}{S_2^2} \sim F(n_1-1, n_2-1) \quad \text{dưới } H_0.
  \]
  Bác bỏ nếu $F_0 > F_{\frac{\alpha}{2}}(n_1-1, n_2-1)$ hoặc $F_0 < F_{1-\frac{\alpha}{2}}(n_1-1, n_2-1)$.
\end{congthuc}

\begin{vidu}[title={So sánh độ phân tán năng suất hai loại cây}]
Loại A: $n_1=15$, $s_1^2=12$. Loại B: $n_2=18$, $s_2^2=25$. $H_0: \sigma_1^2=\sigma_2^2$ đối với $H_1: \sigma_1^2 \ne \sigma_2^2$, $\alpha=5\%$:
\[
  F_0 = \dfrac{12}{25} = 0{,}48.
\]
$F_{0{,}025;14,17} \approx 2{,}86$, $F_{0{,}975;14,17} \approx 0{,}38$. Vì $0{,}38 < 0{,}48 < 2{,}86$, \textbf{chưa bác bỏ} $H_0$.
\end{vidu}

\begin{luuy}[title={Thứ tự F}]
Luôn đặt $S_1^2$ (phương sai lớn hơn) ở tử số để $F_0 \ge 1$, hoặc dùng miền bác bỏ hai phía đối xứng.
\end{luuy}

\subsubsection{So sánh hai tỷ lệ}

\begin{congthuc}[title={Kiểm định hai tỷ lệ - mẫu lớn}]
  $H_0: p_1 = p_2$ đối với $H_1: p_1 \ne p_2$:
  \[
    Z_0 = \dfrac{\hat{p}_1 - \hat{p}_2}{\sqrt{\hat{p}(1-\hat{p})(\dfrac{1}{n_1} + \dfrac{1}{n_2})}}, \quad
    \hat{p} = \dfrac{x_1 + x_2}{n_1 + n_2}.
  \]
  Dưới $H_0$: $Z_0 \approx \mathcal{N}(0,1)$.
\end{congthuc}

\begin{vidu}[title={So sánh tỷ lệ khỏi bệnh hai thuốc}]
Thuốc A: $n_1=200$, $x_1=160$ khỏi ($\hat{p}_1=0{,}80$). Thuốc B: $n_2=180$, $x_2=144$ ($\hat{p}_2=0{,}80$).
$H_0: p_1=p_2$ đối với $H_1: p_1 \ne p_2$, $\alpha=5\%$:
\[
  \hat{p} = \dfrac{304}{380} = 0{,}8, \quad
  z_0 = \dfrac{0{,}80-0{,}80}{\sqrt{0{,}8\cdot 0{,}2(\dfrac{1}{200}+\dfrac{1}{180})}} = 0.
\]
$z_0 = 0 \notin W_{0{,}05}$ $\Rightarrow$ \textbf{chưa bác bỏ} $H_0$: chưa thấy khác biệt tỷ lệ khỏi.
\end{vidu}

\begin{vidu}[title={So sánh tỷ lệ ủng hộ hai khu vực}]
Khu A: $n_1=500$, $\hat{p}_1=0{,}45$. Khu B: $n_2=400$, $\hat{p}_2=0{,}52$:
\[
  \hat{p} = \dfrac{225+208}{900} = 0{,}481, \quad
  z_0 = \dfrac{0{,}45-0{,}52}{\sqrt{0{,}481\cdot 0{,}519(\dfrac{1}{500}+\dfrac{1}{400})}} \approx -2{,}15.
\]
$|z_0| > 1{,}96$ $\Rightarrow$ \textbf{bác bỏ} $H_0$.
\end{vidu}

\begin{congthuc}[title={KTC cho $\mu_1 - \mu_2$ và $p_1 - p_2$}]
  \textbf{Hiệu trung bình} ($\sigma$ biết):
  \[
    (\bar{X}_1-\bar{X}_2) \pm z_{\frac{\alpha}{2}}\sqrt{\dfrac{\sigma_1^2}{n_1} + \dfrac{\sigma_2^2}{n_2}}.
  \]
  \textbf{Hiệu tỷ lệ}:
  \[
    (\hat{p}_1-\hat{p}_2) \pm z_{\frac{\alpha}{2}}\sqrt{\dfrac{\hat{p}_1(1-\hat{p}_1)}{n_1} + \dfrac{\hat{p}_2(1-\hat{p}_2)}{n_2}}.
  \]
\end{congthuc}

\begin{phuongphap}[title={Chọn kiểm định hai mẫu}]
\begin{enumerate}
  \item Một nhóm / hai điều kiện trên cùng đối tượng $\Rightarrow$ \textbf{kiểm định $t$ trên hiệu cặp} (mẫu cặp).
  \item Hai nhóm độc lập, $\sigma_1^2=\sigma_2^2$ $\Rightarrow$ \textbf{kiểm định $t$ gộp phương sai}.
  \item Hai nhóm độc lập, $\sigma_1^2 \ne \sigma_2^2$ $\Rightarrow$ \textbf{kiểm định $t$ Welch} hoặc kiểm định $F$ trước.
  \item Tỷ lệ, $n$ lớn $\Rightarrow$ \textbf{kiểm định $z$ so sánh hai tỷ lệ}.
\end{enumerate}
\end{phuongphap}

\begin{tongket}[title={Tổng kết}]
\begin{enumerate}
  \item Hai mẫu độc lập: $Z_0$ hoặc $T_0$ cho $\mu_1-\mu_2$.
  \item Cặp đôi: hiệu $D_i$, kiểm định $t(n-1)$ trên $\bar{D}$.
  \item Kiểm định $F$ so sánh $\sigma_1^2$, $\sigma_2^2$.
  \item Kiểm định $z$ so sánh $p_1$, $p_2$.
\end{enumerate}
\end{tongket}

\end{document}
